\chapter{1985年上海卷（理）}

\section{填空题}

\begin{problem}

不等式 \(\left|x+2\right|<1\) 的解集是\fillinblank{}.

\begin{answer}
\(\left(-3,-1\right)\)
\end{answer}
\end{problem}

\begin{solution}
由绝对值不等式得 \(-1<x+2<1\)，所以 \(-3<x<-1\). 因此解集为 \(\left(-3,-1\right)\).
\end{solution}

\begin{problem}

函数 \(y=\sqrt[3]{x}-1\) 的反函数是\fillinblank{}.

\begin{answer}
\(y=\left(x+1\right)^3\)
\end{answer}
\end{problem}

\begin{solution}
令原函数中的自变量与函数值互换，得 \(x=\sqrt[3]{y}-1\). 移项并立方，得 \(\sqrt[3]{y}=x+1\)，所以反函数为 \(y=\left(x+1\right)^3\).
\end{solution}

\begin{problem}

点 \((0,1)\) 到直线 \(x+y=2\) 的距离是\fillinblank{}.

\begin{answer}
\(\dfrac{\sqrt{2}}{2}\)
\end{answer}
\end{problem}

\begin{solution}
将直线写成 \(x+y-2=0\). 由点到直线的距离公式，所求距离为
\(\dfrac{\left|0+1-2\right|}{\sqrt{1^2+1^2}}=\dfrac{1}{\sqrt{2}}=\dfrac{\sqrt{2}}{2}\).
\end{solution}

\begin{problem}

函数 \(y=\sin\dfrac{\pi x}{2}\) 的最小正周期是\fillinblank{}.

\begin{answer}
\(4\)
\end{answer}
\end{problem}

\begin{solution}
正弦函数 \(\sin u\) 的最小正周期为 \(2\pi\). 令 \(u=\dfrac{\pi x}{2}\)，当 \(x\) 增加 \(T\) 时，须有 \(\dfrac{\pi T}{2}=2\pi\)，从而 \(T=4\).
\end{solution}

\begin{problem}

如果一个圆的圆心在点 \((2,4)\)，并且经过点 \((0,3)\)，那么这个圆的方程是\fillinblank{}.

\begin{answer}
\(\left(x-2\right)^2+\left(y-4\right)^2=5\)
\end{answer}
\end{problem}

\begin{solution}
半径的平方等于圆心到已知点的距离平方，因此
\(r^2=\left(0-2\right)^2+\left(3-4\right)^2=4+1=5\). 代入圆的标准方程，得 \(\left(x-2\right)^2+\left(y-4\right)^2=5\).
\end{solution}

\begin{problem}

若平面 \(\alpha\) 及这个平面外的一条直线 \(l\) 同时垂直于直线 \(m\)，则直线 \(l\) 和平面 \(\alpha\) 的位置关系是\fillinblank{}.

\begin{answer}
\(l\parallel\alpha\)
\end{answer}
\end{problem}

\begin{solution}
因为平面 \(\alpha\perp m\)，所以直线 \(m\) 是平面 \(\alpha\) 的法线方向. 又因为 \(l\perp m\)，直线 \(l\) 的方向与平面 \(\alpha\) 平行. 题设说明 \(l\) 在平面 \(\alpha\) 外，故 \(l\) 不在该平面内，从而 \(l\parallel\alpha\).
\end{solution}

\begin{problem}

函数 \(y=\sin^{2}x-\cos^{2}x\) 的最小值是\fillinblank{}.

\begin{answer}
\(-1\)
\end{answer}
\end{problem}

\begin{solution}
由二倍角公式，\(\sin^{2}x-\cos^{2}x=-\cos 2x\). 因为 \(-1\leqslant\cos 2x\leqslant1\)，所以 \(-\cos 2x\) 的最小值为 \(-1\)，且在 \(\cos 2x=1\) 时取得.
\end{solution}

\begin{problem}

方程 \(9^{x}-7\cdot3^{x}-18=0\) 的解是\fillinblank{}.

\begin{answer}
\(x=2\)
\end{answer}
\end{problem}

\begin{solution}
令 \(t=3^x\)，则 \(t>0\)，且 \(9^x=t^2\). 原方程化为 \(t^2-7t-18=0\)，即 \((t-9)(t+2)=0\). 由于 \(t>0\)，只能取 \(t=9=3^2\)，所以 \(x=2\).
\end{solution}

\begin{problem}

若 \(\tan\alpha=-2\)，且 \(\sin\alpha<0\)，则 \(\cos\alpha\) 的值是\fillinblank{}.

\begin{answer}
\(\dfrac{\sqrt{5}}{5}\)
\end{answer}
\end{problem}

\begin{solution}
由 \(\tan\alpha<0\) 且 \(\sin\alpha<0\)，可知 \(\cos\alpha>0\). 又
\(1+\tan^2\alpha=\dfrac{1}{\cos^2\alpha}\)，所以 \(\cos^2\alpha=\dfrac{1}{1+4}=\dfrac{1}{5}\). 结合 \(\cos\alpha>0\)，得 \(\cos\alpha=\dfrac{\sqrt{5}}{5}\).
\end{solution}

\begin{problem}

已知 \(O\) 为直角坐标系的原点，点 \(A\) 的坐标为 \((1,3)\)，点 \(B\) 的坐标为 \((0,3)\). 若把 \(\triangle OAB\) 绕 \(y\) 轴旋转一周，则所得旋转体的体积是\fillinblank{}.

\begin{answer}
\(\pi\)
\end{answer}
\end{problem}

\begin{solution}
三角形 \(OAB\) 的边 \(OB\) 在旋转轴 \(y\) 轴上，边 \(AB\) 垂直于旋转轴，且 \(AB=1\)，\(OB=3\). 旋转后所得立体是底面半径为 \(1\)、高为 \(3\) 的圆锥，因此体积为 \(\dfrac{1}{3}\pi\cdot1^2\cdot3=\pi\).
\end{solution}

\section{单选题}

\begin{problem}

设 \(z\) 为复数，\(z\) 是 \(\bar{z}\) 的共轭复数，则 \(z=\bar{z}\) 是 \(z\) 为实数的（\quad）

\choices
    {充分但不必要的条件}
    {必要但不充分的条件}
    {充分而且必要的条件}
    {既不充分也不必要的条件}

\begin{answer}
C
\end{answer}
\end{problem}

\begin{solution}
设 \(z=a+bi\)，其中 \(a,b\) 为实数，则 \(\bar{z}=a-bi\). 因而 \(z=\bar{z}\) 当且仅当 \(a+bi=a-bi\)，即 \(b=0\). 这又当且仅当 \(z=a\) 为实数，所以该条件是充分且必要的，选 C.
\end{solution}

\begin{problem}

下列各式中，正确的是（\quad）

\choices
    {\(\arcsin\left(-\dfrac{\pi}{3}\right)=-\dfrac{\sqrt{3}}{2}\)}
    {\(\sin\left(\arcsin\dfrac{\pi}{3}\right)=\dfrac{\pi}{3}\)}
    {\(\arcsin\left(\sin\dfrac{5\pi}{4}\right)=-\dfrac{\pi}{4}\)}
    {\(\sin\left[\arccos\left(-\dfrac{1}{2}\right)\right]=-\dfrac{\sqrt{3}}{2}\)}

\begin{answer}
C
\end{answer}
\end{problem}

\begin{solution}
反正弦函数的定义域为 \([-1,1]\)，所以 A 中的 \(-\dfrac{\pi}{3}\) 和 B 中的 \(\dfrac{\pi}{3}\) 都不在定义域内. 又 \(\sin\dfrac{5\pi}{4}=-\dfrac{\sqrt{2}}{2}\)，且反正弦函数的值域为 \(\left[-\dfrac{\pi}{2},\dfrac{\pi}{2}\right]\)，故 C 的左边等于 \(-\dfrac{\pi}{4}\). 最后，\(\arccos\left(-\dfrac{1}{2}\right)=\dfrac{2\pi}{3}\)，所以 D 的左边为 \(\dfrac{\sqrt{3}}{2}\). 因此只有 C 正确.
\end{solution}

\begin{problem}

已知函数 \(f(x)=\lg\left(x^2-3x+2\right)\) 的定义域为 \(F\)，函数 \(g(x)=\lg(x-1)+\lg(x-2)\) 的定义域为 \(G\)，那么（\quad）

\choices
    {\(F\cap G=\varnothing\)}
    {\(F=G\)}
    {\(F\subset G\)}
    {\(G\subset F\)}

\begin{answer}
D
\end{answer}
\end{problem}

\begin{solution}
由对数真数须为正，\(f\) 的定义域满足 \(x^2-3x+2=(x-1)(x-2)>0\)，所以 \(F=\left(-\infty,1\right)\cup\left(2,+\infty\right)\). 函数 \(g\) 的两个对数须分别有意义，故 \(x-1>0\) 且 \(x-2>0\)，从而 \(G=\left(2,+\infty\right)\). 因此 \(G\subset F\)，选 D.
\end{solution}

\begin{problem}

若一个棱锥的底面是边数大于 \(3\) 的凸多边形，它的顶点到底面各边的距离都相等，则这个棱锥的底面多边形（\quad）

\choices
    {必为正多边形}
    {必有内切圆}
    {必有外接圆}
    {必既有内切圆又有外接圆}

\begin{answer}
B
\end{answer}
\end{problem}

\begin{solution}
设棱锥顶点为 \(P\)，其在底面上的射影为 \(H\)，底面各边所在直线为 \(e_i\). 设从 \(H\) 到 \(e_i\) 的距离为 \(d_i\)，棱锥的高为 \(h\). 由于 \(PH\) 垂直底面，在垂直于 \(e_i\) 的截面内有 \(PE_i^2=h^2+d_i^2\)，其中 \(E_i\) 是相应的垂足. 题设各个 \(PE_i\) 相等，故各个 \(d_i\) 相等. 凸多边形的内点 \(H\) 到各边距离相等，说明以 \(H\) 为圆心、公共距离为半径的圆同时与各边相切. 所以底面多边形必有内切圆，选 B.
\end{solution}

\begin{problem}

若平移坐标轴，把坐标系 \(xoy\) 的原点 \(O\) 移到点 \(O'\)，\(O'\) 在原坐标系中的坐标为 \((2,-1)\)，则原坐标系中的曲线 \(y=x^3\) 在新坐标系 \(x'o'y'\) 中的方程是（\quad）

\choices
    {\(y'+1=\left(x'-2\right)^3\)}
    {\(y'+1=\left(x'+2\right)^3\)}
    {\(y'-1=\left(x'-2\right)^3\)}
    {\(y'-1=\left(x'+2\right)^3\)}

\begin{answer}
D
\end{answer}
\end{problem}

\begin{solution}
平移后两坐标系的轴方向相同，而新原点在旧坐标系中的坐标为 \((2,-1)\)，所以同一点的坐标满足 \(x=x'+2\)，\(y=y'-1\). 代入原方程 \(y=x^3\)，得 \(y'-1=\left(x'+2\right)^3\)，因此选 D.
\end{solution}

\section{解答题}

\begin{problem}

\begin{enumerate}
\item 在直角坐标系内，方程 \(\begin{cases}
x=3\cot\varphi\\
y=2\csc\varphi
\end{cases}\)（\(\varphi\) 是参数）表示什么曲线？画出它的图形.
\item 在极坐标系内，方程 \(\dfrac{2\sin\theta}{\rho}=1\) 表示什么曲线？画出它的图形.
\end{enumerate}
\end{problem}

\begin{solution}
\begin{enumerate}
\item 由 \(\cot\varphi=\dfrac{x}{3}\)，\(\csc\varphi=\dfrac{y}{2}\)，以及恒等式 \(\csc^2\varphi-\cot^2\varphi=1\)，得 \(\dfrac{y^2}{4}-\dfrac{x^2}{9}=1\). 因而曲线是以原点为中心、实轴在 \(y\) 轴上的双曲线，顶点为 \((0,2)\)，\((0,-2)\)，渐近线为 \(y=\dfrac{2}{3}x\) 和 \(y=-\dfrac{2}{3}x\). 参数取遍允许值时，两支均能得到.
\item 由 \(\dfrac{2\sin\theta}{\rho}=1\)，得 \(\rho=2\sin\theta\). 两边乘以 \(\rho\)，并利用 \(y=\rho\sin\theta\) 和 \(\rho^2=x^2+y^2\)，得到 \(x^2+y^2=2y\)，即 \(x^2+\left(y-1\right)^2=1\). 因而曲线是圆心为 \((0,1)\)、半径为 \(1\) 的圆.
\end{enumerate}
\end{solution}

\begin{problem}

\begin{enumerate}
\item 求 \(\left(x^{3}-\dfrac{1}{x^{2}}\right)^{15}\) 的展开式中的常数项.

\item 从六个数字 \(1\)，\(2\)，\(3\)，\(4\)，\(5\)，\(6\) 中任取四个不同的数字，有多少种取法？由这六个数字可以组成多少个没有重复数字的四位偶数？

\item 在一列球中，第 1 个球的半径为 \(1\)，第 2 个球的直径是第 1 个球的半径，以后每一个球的直径都是前一个球的半径，这样无限继续下去，求所有这些球的表面积的和.
\end{enumerate}
\end{problem}

\begin{solution}
\begin{enumerate}
\item 展开式通项中，若取 \(k\) 个 \(-\dfrac{1}{x^2}\)，则该项为 \(\mathrm C_{15}^{k}\left(x^3\right)^{15-k}\left(-\dfrac{1}{x^2}\right)^k\)，其中 \(x\) 的指数为 \(3(15-k)-2k=45-5k\). 要成为常数项，须有 \(45-5k=0\)，即 \(k=9\). 所以常数项为 \((-1)^9\mathrm C_{15}^{9}=-\mathrm C_{15}^{6}=-5005\).
\item 任取四个不同的数字，取法数为 \(\mathrm C_6^4=15\). 四位偶数的个位必须从 \(2\)，\(4\)，\(6\) 中选取，有 \(3\) 种；个位确定后，千位、百位、十位从余下的五个数字中依次选取，排列数为 \(5\cdot4\cdot3\). 因此四位偶数共有 \(3\cdot5\cdot4\cdot3=180\) 个.
\item 设第 \(n\) 个球的半径为 \(r_n\). 由题意，\(r_1=1\)，且 \(r_{n+1}=\dfrac{r_n}{2}\)，所以 \(r_n=\dfrac{1}{2^{n-1}}\). 所有球的表面积之和为 \(4\pi\sum_{n=1}^{\infty}r_n^2=4\pi\sum_{n=1}^{\infty}\left(\dfrac14\right)^{n-1}=\dfrac{4\pi}{1-\frac14}=\dfrac{16\pi}{3}\).
\end{enumerate}
\end{solution}

\begin{problem}

设 \(\triangle ABC\) 的两个内角 \(A\)，\(B\) 所对的边分别为 \(a\)，\(b\). 若复数 \(z_1\) 的模为 \(a\)，幅角为 \(B\)；复数 \(z_2\) 的模为 \(b\)，幅角为 \(-A\). 求证 \(z_1+z_2\) 为实数.
\end{problem}

\begin{solution}
由复数的三角形式，\(z_1=a\left(\cos B+\i\sin B\right)\)，\(z_2=b\left(\cos A-\i\sin A\right)\). 因此
\(z_1+z_2=\left(a\cos B+b\cos A\right)+\i\left(a\sin B-b\sin A\right)\).

由正弦定理，存在外接圆半径 \(R\)，使得 \(a=2R\sin A\)，\(b=2R\sin B\). 所以 \(a\sin B=b\sin A=2R\sin A\sin B\)，从而 \(z_1+z_2\) 的虚部为零，故 \(z_1+z_2\) 为实数.
\end{solution}

\begin{problem}

解三角方程 \(\cos^2x+\cos^23x=1\).
\end{problem}

\begin{solution}
由半角变换，原方程等价于 \(\dfrac{1+\cos2x}{2}+\dfrac{1+\cos6x}{2}=1\)，即 \(\cos2x+\cos6x=0\). 利用和差化积，得 \(2\cos4x\cos2x=0\). 因而
\(\cos4x=0\) 或 \(\cos2x=0\)，分别得到 \(x=\dfrac{\pi}{8}+\dfrac{k\pi}{4}\) 或 \(x=\dfrac{\pi}{4}+\dfrac{k\pi}{2}\)，其中 \(k\in\Z\). 所以原方程的解为这两类解的并集.
\end{solution}

\begin{problem}

已知函数 \(f(x)\) 对其定义域内的任意两个实数 \(a\)，\(b\)，当 \(a<b\) 时，都有 \(f(a)<f(b)\). 试用反证法证明：方程 \(f(x)=0\) 至多有一个实数根.
\end{problem}

\begin{solution}
假设方程 \(f(x)=0\) 至少有两个不同的实数根，设它们为 \(x_1\)，\(x_2\). 不妨设 \(x_1<x_2\). 由题设函数严格递增，应有 \(f(x_1)<f(x_2)\). 但 \(x_1\)，\(x_2\) 都是方程的根，所以 \(f(x_1)=f(x_2)=0\)，这与 \(f(x_1)<f(x_2)\) 矛盾. 因此假设不成立，方程 \(f(x)=0\) 至多有一个实数根.
\end{solution}

\begin{problem}

对于一切大于 \(1\) 的自然数 \(n\)，证明：\(\left(1+\dfrac{1}{3}\right)\left(1+\dfrac{1}{5}\right)\cdots\left(1+\dfrac{1}{2n-1}\right)>\dfrac{\sqrt{2n+1}}{2}\).
\end{problem}

\begin{solution}
记 \(P_n=\displaystyle\prod_{j=2}^{n}\left(1+\dfrac{1}{2j-1}\right)=\prod_{j=2}^{n}\dfrac{2j}{2j-1}\). 当 \(n=2\) 时，\(P_2=\dfrac43>\dfrac{\sqrt5}{2}\)，因为 \(\dfrac{16}{9}>\dfrac54\).

假设对某个大于等于 \(2\) 的自然数 \(n\) 有 \(P_n>\dfrac{\sqrt{2n+1}}{2}\). 注意到
\(\dfrac{P_{n+1}}{P_n}=\dfrac{2n+2}{2n+1}\)，且
\(\left(\dfrac{2n+2}{2n+1}\right)^2-\dfrac{2n+3}{2n+1}=\dfrac{1}{(2n+1)^2}>0\). 因而 \(\dfrac{2n+2}{2n+1}>\sqrt{\dfrac{2n+3}{2n+1}}\). 结合归纳假设，得
\(P_{n+1}>\dfrac{\sqrt{2n+1}}{2}\sqrt{\dfrac{2n+3}{2n+1}}=\dfrac{\sqrt{2n+3}}{2}\).

由数学归纳法，所证不等式对一切大于 \(1\) 的自然数 \(n\) 成立.
\end{solution}

\begin{problem}

在三棱锥 \(P-ABC\) 中，顶点 \(P\) 到 \(BC\)，\(AC\)，\(AB\) 的距离分别为 \(h_1\)，\(h_2\)，\(h_3\)，二面角 \(P-BC-A=\alpha_1\)，\(P-AC-B=\alpha_2\)，\(P-AB-C=\alpha_3\). \(\alpha_1\)，\(\alpha_2\)，\(\alpha_3\) 都是锐角，且依次成等差数列，\(h_1\)，\(h_2\)，\(h_3\) 依次成等比数列. 求证：\(\alpha_1=\alpha_2=\alpha_3\).
\end{problem}

\begin{solution}
记三棱锥体积为 \(V\)，底面三角形面积为 \(S\). 设从 \(A\) 到 \(BC\) 的高为 \(h_a\). 以 \(BC\) 为棱作垂直截面，则点 \(P\) 到 \(BC\) 的距离为 \(h_1\)，而点 \(A\) 到平面 \(PBC\) 的距离为 \(h_a\sin\alpha_1\). 所以
\(V=\dfrac13S_{\triangle PBC}h_a\sin\alpha_1=\dfrac13S h_1\sin\alpha_1\).

同理，\(V=\dfrac13S h_2\sin\alpha_2=\dfrac13S h_3\sin\alpha_3\)，从而
\(h_1\sin\alpha_1=h_2\sin\alpha_2=h_3\sin\alpha_3\). 由于 \(h_1\)，\(h_2\)，\(h_3\) 成等比数列，\(h_2^2=h_1h_3\). 设上述公共值为 \(c\)，则
\(\sin^2\alpha_2=\dfrac{c^2}{h_2^2}=\dfrac{c^2}{h_1h_3}=\sin\alpha_1\sin\alpha_3\).

又因 \(\alpha_1\)，\(\alpha_2\)，\(\alpha_3\) 成等差数列，\(\alpha_2=\dfrac{\alpha_1+\alpha_3}{2}\). 令 \(d=\dfrac{\alpha_3-\alpha_1}{2}\)，则
\(\sin\alpha_1\sin\alpha_3=\sin\left(\alpha_2-d\right)\sin\left(\alpha_2+d\right)=\sin^2\alpha_2-\sin^2d\). 与前式比较，得 \(\sin^2d=0\). 因为两角均为锐角，故 \(d=0\)，即 \(\alpha_1=\alpha_3\). 再由等差关系，得到 \(\alpha_1=\alpha_2=\alpha_3\).
\end{solution}

\begin{problem}

已知直线 \(y=x+m\) 和曲线 \(x^3+2y^3+y^4-1=0\) 交于 \(A\)，\(B\) 两点，\(P\) 是这条直线上的点，且 \(\left|PA\right|\cdot\left|PB\right|=2\). 求：当 \(m\) 变化时，点 \(P\) 的轨迹方程，并说明轨迹是什么图形.
\end{problem}

\begin{solution}
设 \(P=(x,y)\). 直线方向向量为 \((1,1)\)，故可设
\(A=(x+u,y+u)\)，\(B=(x+v,y+v)\). 于是 \(\left|PA\right|=\sqrt{2}\left|u\right|\)，\(\left|PB\right|=\sqrt{2}\left|v\right|\)，已知条件等价于 \(\left|uv\right|=1\). 记 \(\varepsilon=uv\)，则 \(\varepsilon=1\) 或 \(\varepsilon=-1\)，并记 \(s=u+v\).

把点 \((x+t,y+t)\) 代入曲线方程，得到
\(G(t)=t^4+c_3t^3+c_2t^2+c_1t+c_0\)，其中
\(c_3=4y+3\)，\(c_2=3x+6y^2+6y\)，\(c_1=3x^2+4y^3+6y^2\)，\(c_0=x^3+2y^3+y^4-1\).

由于 \(t=u\)，\(t=v\) 分别对应交点 \(A\)，\(B\)，且 \(uv=\varepsilon\)，所以二次式 \(t^2-st+\varepsilon\) 是 \(G(t)\) 的因式. 设另一因式为 \(t^2+qt+r\)，比较常数项和三次项，得 \(q=c_3+s\)，\(r=\varepsilon c_0\)，从而
\(G(t)=\left(t^2-st+\varepsilon\right)\left(t^2+(c_3+s)t+\varepsilon c_0\right)\).

比较二次项和一次项，得到关于 \(s\) 的两个条件
\[
\begin{cases}
s^2+c_3s+c_2-\varepsilon c_0-\varepsilon=0,\\
\varepsilon(1-c_0)s+\varepsilon c_3-c_1=0.
\end{cases}
\]
这就是轨迹的紧凑方程表示，其中 \(\varepsilon=1\) 和 \(\varepsilon=-1\) 分别对应点 \(P\) 在交点 \(A\)，\(B\) 的同侧和异侧.

当 \(c_0\neq1\) 时，由第二式消去 \(s\)，得
\(s=\dfrac{\varepsilon c_1-c_3}{1-c_0}\). 代回第一式，轨迹方程可写为
\[
\left(\varepsilon c_1-c_3\right)^2+c_3\left(\varepsilon c_1-c_3\right)\left(1-c_0\right)+\left(c_2-\varepsilon c_0-\varepsilon\right)\left(1-c_0\right)^2=0,
\qquad \varepsilon=\pm1,
\]
其中 \(c_0,c_1,c_2,c_3\) 如上定义；当 \(c_0=1\) 时使用前面的二元条件组. 同时还须取使 \(t^2-st+\varepsilon\) 有两个不同实根、而另一因式没有实根的分支，即
\(s^2-4\varepsilon>0\)，\((c_3+s)^2-4\varepsilon c_0<0\).

因此，点 \(P\) 的轨迹是上述两种符号对应的十二次代数曲线的实分支，并按最后两个不等式取出与题设“直线交于 \(A\)，\(B\) 两点”相符的分支；它不是圆、椭圆、抛物线或双曲线等二次曲线.
\end{solution}
